\documentclass[11pt,a4paper]{article}
\usepackage[T1]{fontenc}
\usepackage[utf8]{inputenc}
\usepackage{amsmath,amssymb,amsthm}
\usepackage[margin=1in]{geometry}
\usepackage[colorlinks=true,linkcolor=blue,urlcolor=blue]{hyperref}
\theoremstyle{plain}
\newtheorem{theorem}{Theorem}
\newtheorem{lemma}[theorem]{Lemma}
\newtheorem{proposition}[theorem]{Proposition}
\newtheorem{corollary}[theorem]{Corollary}
\theoremstyle{definition}
\newtheorem{remark}[theorem]{Remark}
% A quoted conjecture can be a single hundred-term formula whose terms are thirty-digit
% numbers. TeX cannot hyphenate those, so without a little stretch every such quote runs off
% the page; the linked-a-file papers are all of that shape.
\setlength{\emergencystretch}{3em}

\title{The empirical closed form for OEIS A223847, proved}
\author{Adrian Perez Fontelles\\ \small Independent researcher}
\date{1 September 2026}
\begin{document}
\maketitle

\begin{abstract}
OEIS A223847 carries an empirical closed form for $a(n)$ --- not a recurrence relating
consecutive terms but an explicit formula. It is true, and it is decidable rather than
empirical. The entry counts arrays subject to a condition on a bounded window of consecutive lines, so the lines of an array are the
vertices of a finite digraph, an array is a walk in it, and $a(n)$ is a walk count on
$S=9472$ vertices. Any function of the form $\sum_j p_j(n)\lambda_j^{\,n}$ satisfies the
monic linear recurrence whose characteristic polynomial is
$\prod_j(x-\lambda_j)^{\deg p_j+1}$; here that polynomial is $q(x)=(x-1)^{31}$, of degree
$31$. Two things are then checked exactly: that the walk count satisfies the same
recurrence from some index on, which the Cayley--Hamilton theorem reduces to a finite
computation, and that the two sequences agree at $31$ consecutive indices past that
point. A monic recurrence determines every later term from its predecessors, so agreement
there is agreement for good.
\end{abstract}

\noindent\small 2020 Mathematics Subject Classification. 05A15, 11B37, 15A18.\normalsize

\section{The conjecture}

OEIS A223847 is ``Number of nX5 0..3 arrays with rows and columns unimodal.''. It has offset $1$ and begins
\[
296, 87616, 10804883, 664770145, 24291817048, 594861333098, 10609537390768, 146233793223364,\ \dots
\]
The entry states, as a conjecture contributed by its author and never marked settled:
\begin{quote}\raggedright\small
Empirical: a(n) = (12185409283/15043832793341144432640000000)*n\^{}30 + (76296730104643/884176199373970195454361600000)*n\^{}29 + (1070597323435537/228666258458785395376128000000)*n\^{}28 + (15095976776879/89620324694801252352000000)*n\^{}27 + (80978639299446029/18148115750697253601280000000)*n\^{}26 + (17128733527118137/186134520519971831808000000)*n\^{}25 + (1793728053298859387/1172647479275822540390400000)*n\^{}24 + (1261636802562641833/60135768167990899507200000)*n\^{}23 + (5594512626040336997/23174851369087401984000000)*n\^{}22 + (31544771499812381/13376537586774835200000)*n\^{}21 + (47886425657641284787/2427841571999632588800000)*n\^{}20 + (229975881467092172089/1618561047999755059200000)*n\^{}19 + (9082589031754852696739/10260073516345196544000000)*n\^{}18 + (111240879091929267206971/23256166637049112166400000)*n\^{}17 + (46080377535895002055433/2052014703269039308800000)*n\^{}16 + (5972995162066161572851/65143323913302835200000)*n\^{}15 + (2705387734319931209037779/8305773798946111488000000)*n\^{}14 + (21438489128749966497727/21296855894733619200000)*n\^{}13 + (123748166982394715917397/45739838228461977600000)*n\^{}12 + (281638905229159746619297/44567021863629619200000)*n\^{}11 + (1531541445561308485678126639/119495327371856916480000000)*n\^{}10 + (569764139394721423500017/25290016374996172800000)*n\^{}9 + (12437437033053107961181039/363543985390569984000000)*n\^{}8 + (1801415998853685830112073/40393776154507776000000)*n\^{}7 + (909519212873314420568666851/18379168150301038080000000)*n\^{}6 + (1564096450989494220884147/34035496574631552000000)*n\^{}5 + (108960453802905313513657/3063194691716839680000)*n\^{}4 + (2031598190834852737/90040996229184000)*n\^{}3 + (318817655518333/25794785016000)*n\^{}2 + (1721120935577/388181593800)*n + 1
\end{quote}
As of the ``Last modified'' line on the live entry (Oct 03 2025 22:43:42, revision 8) this is still
recorded as empirical, and nothing on the entry records it as proved. Write
\begin{multline*}
f(n)\;=\;\frac{12185409283 n^{30}}{15043832793341144432640000000} + \frac{76296730104643 n^{29}}{884176199373970195454361600000} \\
+ \frac{1070597323435537 n^{28}}{228666258458785395376128000000} + \frac{15095976776879 n^{27}}{89620324694801252352000000} \\
+ \frac{80978639299446029 n^{26}}{18148115750697253601280000000} + \frac{17128733527118137 n^{25}}{186134520519971831808000000} \\
+ \frac{1793728053298859387 n^{24}}{1172647479275822540390400000} + \frac{1261636802562641833 n^{23}}{60135768167990899507200000} \\
+ \frac{5594512626040336997 n^{22}}{23174851369087401984000000} + \frac{31544771499812381 n^{21}}{13376537586774835200000} \\
+ \frac{47886425657641284787 n^{20}}{2427841571999632588800000} + \frac{229975881467092172089 n^{19}}{1618561047999755059200000} \\
+ \frac{9082589031754852696739 n^{18}}{10260073516345196544000000} + \frac{111240879091929267206971 n^{17}}{23256166637049112166400000} \\
+ \frac{46080377535895002055433 n^{16}}{2052014703269039308800000} + \frac{5972995162066161572851 n^{15}}{65143323913302835200000} \\
+ \frac{2705387734319931209037779 n^{14}}{8305773798946111488000000} + \frac{21438489128749966497727 n^{13}}{21296855894733619200000} \\
+ \frac{123748166982394715917397 n^{12}}{45739838228461977600000} + \frac{281638905229159746619297 n^{11}}{44567021863629619200000} \\
+ \frac{1531541445561308485678126639 n^{10}}{119495327371856916480000000} + \frac{569764139394721423500017 n^{9}}{25290016374996172800000} \\
+ \frac{12437437033053107961181039 n^{8}}{363543985390569984000000} + \frac{1801415998853685830112073 n^{7}}{40393776154507776000000} \\
+ \frac{909519212873314420568666851 n^{6}}{18379168150301038080000000} + \frac{1564096450989494220884147 n^{5}}{34035496574631552000000} \\
+ \frac{108960453802905313513657 n^{4}}{3063194691716839680000} + \frac{2031598190834852737 n^{3}}{90040996229184000} \\
+ \frac{318817655518333 n^{2}}{25794785016000} + \frac{1721120935577 n}{388181593800} \\
+ 1.
\end{multline*}

\section{The count is a walk count}

The entry's defining condition is local: it is a condition on a bounded window of consecutive lines. Take as vertices the
admissible configurations of such a window and put an edge from $u$ to $v$ when $v$ may
follow $u$, which the condition decides by inspection. An admissible array is then exactly a
walk, so with $M$ the adjacency matrix of that digraph on $S=9472$ vertices, and $u,w$ the
vectors recording which windows may start and end an array,
\[
a(n)\;=\;u^{\!\top}M^{\,g(n)}w
\]
for the linear function $g$ that the entry's shape supplies. In particular $a$ is
$C$-finite: it satisfies some monic integer linear recurrence, though not necessarily a short
one.

\section{A closed form is a recurrence}

\begin{lemma}\label{lem:ann}
Let $f(m)=\sum_{j}p_j(m)\lambda_j^{\,m}$ with the $\lambda_j$ distinct and the $p_j$
polynomials, and put $q(x)=\prod_j(x-\lambda_j)^{\deg p_j+1}$. Then $q$ applied to the shift
operator annihilates $f$: writing $q(x)=x^{r}-\sum_{i=1}^{r}c_ix^{r-i}$,
\[
f(m)\;=\;\sum_{i=1}^{r}c_i\,f(m-i)\qquad\text{for every }m.
\]
\end{lemma}

\begin{proof}
The shift operator $E$ acts on $m\mapsto\lambda^{m}p(m)$ as
$(E-\lambda)\bigl(\lambda^{m}p(m)\bigr)=\lambda^{m+1}\bigl(p(m+1)-p(m)\bigr)$, and
$p(m+1)-p(m)$ has degree one less than $p$. So $(E-\lambda)^{\deg p+1}$ kills
$\lambda^{m}p(m)$, and the factors of $q$ commute.
\end{proof}

Here $q(x)=(x-1)^{31}$, of degree $31$; the closed form is a polynomial in $n$, so this is $(x-1)^{31}$, and the recurrence it encodes is
\begin{equation}\label{eq:rec}
a(n)\;=\;31\,a(n-1) - 465\,a(n-2) + 4495\,a(n-3) - 31465\,a(n-4) + 169911\,a(n-5) - 736281\,a(n-6) + 2629575\,a(n-7) - 7888725\,a(n-8) + 20160075\,a(n-9) - 44352165\,a(n-10) + 84672315\,a(n-11) - 141120525\,a(n-12) + 206253075\,a(n-13) - 265182525\,a(n-14) + 300540195\,a(n-15) - 300540195\,a(n-16) + 265182525\,a(n-17) - 206253075\,a(n-18) + 141120525\,a(n-19) - 84672315\,a(n-20) + 44352165\,a(n-21) - 20160075\,a(n-22) + 7888725\,a(n-23) - 2629575\,a(n-24) + 736281\,a(n-25) - 169911\,a(n-26) + 31465\,a(n-27) - 4495\,a(n-28) + 465\,a(n-29) - 31\,a(n-30) + a(n-31).
\end{equation}

\section{The proof}

\begin{lemma}\label{lem:crit}
With $q$ as above, $a$ satisfies \eqref{eq:rec} for all large $n$ if and only if
$u^{\!\top}M^{\,j}q(M)w=0$ for every $j\ge0$, and it is enough to check $j<S$.
\end{lemma}

\begin{proof}
Substituting the walk expression, the residual $a(n)-\sum_ic_ia(n-i)$ equals
$u^{\!\top}M^{\,g(n)-r}q(M)w$ up to the fixed linear reparametrisation $g$. For the bound,
$M^{S}$ is an integer combination of $I,M,\dots,M^{S-1}$ by Cayley--Hamilton, so the values
$u^{\!\top}M^{j}q(M)w$ satisfy the monic recurrence given by the characteristic polynomial of
$M$; $S$ consecutive zeros therefore force all later ones, and the last nonzero value pins the
threshold exactly.
\end{proof}

\begin{theorem}
$a(n)=f(n)$ for every $n\ge1$.
\end{theorem}

\begin{proof}
The vector $w'=q(M)w$ was formed by $31$ matrix--vector products in exact integer
arithmetic, and $u^{\!\top}M^{j}w'$ evaluated for $j=0,1,\dots$, again exactly. Those integers
vanish from the point corresponding to $n=32$ onwards, and the run was continued until the
working vector was identically zero, which settles every larger $j$ at once. So $a$ satisfies
\eqref{eq:rec} for every $n>31$. By Lemma~\ref{lem:ann} $f$ satisfies \eqref{eq:rec}
identically. Both were then evaluated in exact arithmetic at the $31$ consecutive indices
$n=32,\dots,62$ and agree at each. A monic recurrence of order $31$
determines $a(m)$ from $a(m-1),\dots,a(m-31)$, and for every $m\ge63$ both
$m>31$ and $m-31\ge32$ hold, so induction on $m$ gives $a(m)=f(m)$ for every
$m\ge32$.
Below $n=32$ the recurrence argument gives nothing, since \eqref{eq:rec} is only
established for $a$ from $n=32$ on. The remaining indices $n=1,\dots,31$ were
therefore checked one at a time, directly against the walk count in exact integer arithmetic,
and agree.
\end{proof}

The range is tight in the sense that was checked: the closed form holds from the entry's first term onwards.

\section{Verification}

Four checks, each independent of the algebra above.

First, the walk model reproduces every one of the $14$ terms the entry publishes,
exactly, in integer arithmetic. This is what ties the model to the entry: the model is built
by reading the entry's English, and a misreading gives a different digraph and different
counts.

Second, the closed form was evaluated in exact rational arithmetic --- not floating point ---
at every published term from $n=1$ on, and agrees with the entry's own DATA at each.

Third, the annihilation test of Lemma~\ref{lem:crit} was carried out exactly, and the model
and the closed form were then compared directly at sixty consecutive indices beyond
$n=1$, far past where the proof already settles the matter.

Fourth, the closed form was deliberately perturbed --- by a constant and by a term linear in
$n$ --- and each perturbed form was rejected by the same comparison, so the test is not one
that anything would pass.

\begin{thebibliography}{9}
\bibitem{oeis} The OEIS Foundation, \emph{The On-Line Encyclopedia of Integer
Sequences}, \url{https://oeis.org}, 2026. Sequence A223847.
\bibitem{stanley} R.~P.~Stanley, \emph{Enumerative Combinatorics}, Volume 1, 2nd ed.,
Cambridge University Press, 2011. (Transfer-matrix method, Section 4.7; $C$-finite sequences,
Section 4.1.)
\bibitem{kp} M.~Kauers and P.~Paule, \emph{The Concrete Tetrahedron}, Springer, 2011.
(Closed forms and linear recurrences with constant coefficients.)
\bibitem{horn} R.~A.~Horn and C.~R.~Johnson, \emph{Matrix Analysis}, 2nd ed., Cambridge
University Press, 2013.
\end{thebibliography}

\end{document}
